\documentclass[10pt]{article}

\usepackage[utf8]{inputenc}
\usepackage[english]{babel}
\usepackage[margin=1in]{geometry}
\usepackage{amsmath,amsthm,amssymb,stmaryrd}
\usepackage{hyperref,enumitem,xcolor}

\hypersetup{colorlinks,linkcolor={red!50!black},citecolor={blue!50!black},urlcolor={blue!80!black}}

% Plain amsthm: thmtools 0.76 fails with the 2026-06-01 LaTeX kernel on shared counters.
\newtheorem{theorem}{Theorem}[section]
\newtheorem{lemma}[theorem]{Lemma}
\newtheorem{corollary}[theorem]{Corollary}
\theoremstyle{definition}
\newtheorem{definition}[theorem]{Definition}
\newtheorem{question}[theorem]{Question}
\theoremstyle{remark}
\newtheorem{remark}[theorem]{Remark}
\setlist[itemize]{parsep=0pt,topsep=2pt}
\setlist[enumerate]{parsep=0pt,topsep=2pt}

\newcommand{\N}{\mathbb{N}}
\newcommand{\Prop}{\mathsf{Prop}}
\newcommand{\Type}{\mathsf{Type}}
\newcommand{\ZF}{\mathrm{ZF}}
\newcommand{\ZFC}{\mathrm{ZFC}}
\newcommand{\EM}{\mathrm{EM}}
\newcommand{\Con}{\mathrm{Con}}
\newcommand{\rk}{\operatorname{rank}}
\newcommand{\Acc}{\mathsf{Acc}}
\newcommand{\V}{\mathbb{V}}
\newcommand{\la}{\mathsf{a}}
\newcommand{\lb}{\mathsf{b}}
\newcommand{\lc}{\mathsf{c}}
\newcommand{\lean}[1]{\texttt{#1}}
\newcommand{\guard}[2]{\{#1\mid #2\}}

\begin{document}

\title{Accessibility proves Replacement:\\
  \large the consistency of $\ZF$ in type theory with excluded middle and no choice}
\author{Mario Carneiro}
% Drafted with Claude; adjust authorship/acknowledgement as you see fit.
\date{Draft, \today}
\maketitle

\begin{abstract}
The sets-as-trees interpretation of set theory in a dependent type
theory with an impredicative universe of propositions validates
Zermelo set theory, and it validates Replacement if the type theory
has a choice or description operator, which turns a functional relation
into a function. It has been natural to expect that without such an
operator the strength of the type theory drops well below that of $\ZF$. We
show that it does not. In the type theory of Lean with two predicative
universes, from excluded middle as the only assumption and with no
axiom (no choice, no propositional extensionality, no quotients), we
prove the consistency of $\ZF$, stated outright for a first-order proof
system. The proof is formalized. The mechanism is the large elimination
of the accessibility predicate over a type as large as the type of
sets: a recursion on accessibility whose recursive calls are guarded by
propositions, and whose later calls are indexed by the \emph{value} of
an earlier call, computes as a term any ordinal that is specified by a
proposition through a well-founded tree of a certain shape. We give a
rule that produces such a tree for every ordinal, unless some $V_\rho$ is
already a model of $\ZF$; the rule does not choose a cofinal map into a
limit ordinal but takes all definable ones at once. In the first case
the sets-as-trees satisfy Replacement for arbitrary propositional
relations. Either way $\ZF$ has a model.
\end{abstract}

\section{Introduction}

\paragraph{The question.}
Consider a dependent type theory with an impredicative, proof-irrelevant
universe $\Prop$, predicative universes $\Type_0:\Type_1:\cdots$,
inductive types, and an accessibility predicate $\Acc$ in $\Prop$ that
eliminates into every universe: the Calculus of Inductive Constructions
in the form implemented by Lean \cite{lean4,carneiro19}. How strong is
it, as a foundation, when excluded middle is assumed but no choice
principle is?

With choice the answer is known. Aczel's interpretation of sets as
well-founded trees \cite{aczel78} gives a type $\V$ of sets in which
the axioms of Zermelo set theory hold, by impredicativity of $\Prop$ for
Separation and Power set. Replacement asks for the image of a set $s$
under a functional relation $\varphi$. A set in $\V$ is a tree
$\mathsf{sup}\,A\,f$ with $A$ a small type and $f:A\to\V$ a
\emph{function}, so to form the image one needs a function $a\mapsto y_a$
with $\varphi(a,y_a)$, and a functional relation in $\Prop$ does not
provide one, because propositions do not eliminate into types. A choice
or description operator closes the gap, and with it the type theory
with $n+2$ universes interprets $\ZFC$ with $n$ inaccessible cardinals
\cite{werner97,barras10,carneiro19}. Without it, the gap between a
functional relation and a function looks like exactly the gap between
Zermelo set theory and $\ZF$, and the standard set-theoretic model, in
which the universes are $V_\kappa$ for inaccessible $\kappa$, does not tell
whether the inaccessibles are needed.

\paragraph{The result.}
They are needed, in the sense that the type theory with two universes
and excluded middle, without choice, already proves the consistency of
$\ZF$.

\begin{theorem}[formalized]\label{thm:main}
In the type theory of Lean~4, using the universes $\Prop$, $\Type_0$ and
$\Type_1$ and no axioms, excluded middle implies $\Con(\ZF)$, where $\ZF$
is the usual first-order theory with the schemas of Separation and
Replacement and $\Con$ is non-derivability of falsity in a Hilbert
system. More precisely, excluded middle implies that one of the
following holds.
\begin{enumerate}
\item Every functional relation $\varphi:\V\to\V\to\Prop$ on a set $s:\V$ has
  an image in $\V$. Then $\V$ is a model of $\ZF$, in fact of its
  second-order form.
\item There is an ordinal $\rho>\omega$ of $\V$, a limit, such that $V_\rho$ is
  closed under the images of functions definable over $V_\rho$ with
  parameters. Then $V_\rho$ is a model of $\ZF$.
\end{enumerate}
\end{theorem}

By G\"odel's second incompleteness theorem, $\ZF$, if consistent, does
not prove the consistency of this type theory, with or without choice,
so no model construction carried out in $\ZF$ can succeed for it.

\paragraph{The mechanism.}
Propositions do not eliminate into types, with one family of exceptions:
inductive propositions with at most one constructor whose arguments are
propositions or are determined by the indices. Accessibility is the
important one. For a relation $R$ on a type $X$, $\Acc_R\,x$ has the one
constructor $(\forall y.\ R\,y\,x\to\Acc_R\,y)\to\Acc_R\,x$, and its eliminator
defines a function on the accessible elements of $X$ with values in any
type, by recursion on $R$. This is what makes well-founded recursion
available, and it is known to extract information from propositions:
the standard library of Coq obtains a witness of $\exists n.\,P\,n$, for
decidable $P$ on $\N$, by recursion on a proof of an inductive
proposition of the same kind \cite{constructiveepsilon}. That is unbounded search, and it needs the
decidability of $P$ to know when to stop.

We use the same idea transfinitely, with the type of sets itself as the
type of values, and two observations replace decidability.
\begin{itemize}
\item \emph{Guards.} For a proposition $P$ and $t:P\to\V$ there is a set
  $\guard{t}{P}$ which is $t$ if $P$ holds and empty otherwise, namely
  the union of the family $t$ indexed by the proofs of $P$. No decision is made:
  the index type of the tree is $P$ itself.
\item \emph{Adaptive calls.} A step of the recursion may call the
  recursion at a child, obtain a set $G$ as the value, and then call it
  at further children indexed by the elements of a set computed from
  $G$. The index types of sets are small, but there is no bound on the
  sets that occur as values, so the trees that such a recursion explores
  are not bounded in advance by anything the recursion starts from.
\end{itemize}
The \emph{shape} of the tree explored by the recursion, that is, which
guards hold, is given by a relation $R$ on paths which is an arbitrary
proposition. If $R$ is the tree of a well-founded specification of an
ordinal $\eta$, in a sense made precise below, then the root is accessible
and the value of the recursion at the root is $\eta$: an ordinal that was
only specified becomes the value of a term
(Lemma~\ref{lem:mat}). Applied to the ranks of the values of a
functional relation on a set $s$, with one child of the root for each
element of $s$, the value at the root bounds all the ranks, and the
image is obtained by Separation (Theorem~\ref{thm:repl}).

What remains is to specify every ordinal by such a tree, uniformly,
because the trees for the different elements of $s$ must be given by
one proposition. At a limit ordinal $\eta$ the children must have targets
cofinal in $\eta$ and be indexed by a set of rank below $\eta$. We cannot
choose a cofinal map. Instead the children are indexed by \emph{all}
triples of a first-order definition with parameters, a domain, and an
argument, such that the definition gives over $V_\eta$ a function on the
domain with values of unbounded rank; the target of a child is the rank
of the value at the argument. This fails at $\eta$ exactly if no definable
function from a set in $V_\eta$ is unbounded in $V_\eta$, and then $V_\eta$ is a
model of $\ZF$ (Section~\ref{sec:rule}).

\paragraph{The carrier must be big.}
The recursion runs on paths whose labels are arbitrary sets, so its
carrier lives in the same universe $\Type_1$ as $\V$. This is essential.
On a carrier $X$ in $\Type_0$ the height of any accessible element is
below the Hartogs number of the power set of $X$, which is an ordinal of
$\V$ obtained by a term without Replacement, so the tree is bounded in
advance.
We do not know the strength of the type theory in which accessibility is
restricted to carriers in the lowest universe
(Question~\ref{q:small}).

\paragraph{Formalization.}
The development is in Lean~4 without Mathlib, about 2000 lines, and
is described in Section~\ref{sec:lean}; it is available at
\url{https://github.com/digama0/ConZF}. The statement of the main
theorem is
\begin{verbatim}
theorem PSet.con_ZF (em : forall p : Prop, p \/ Not p) : Con ZF
\end{verbatim}
and Lean reports that it depends on no axioms. Everything in
Sections~\ref{sec:sets} to~\ref{sec:zf} is formalized; the text follows
the formal development and names the corresponding declarations.

\section{Sets as trees}\label{sec:sets}

\paragraph{The type theory.}
We use $\Prop$ with its impredicative quantifier, two predicative
universes, inductive types with structural recursion, and $\Acc$ with
elimination into $\Type_1$. The inductive types used are natural numbers,
lists, subtypes, the lifting of a proposition or a small type to a
universe, a three-constructor type of labels, the syntax of formulas,
and the type $\V$ below; all of them are instances of $W$-types and
finite sums. Excluded middle is a hypothesis
$\mathit{em}:\forall p{:}\Prop.\ p\vee\neg p$ of the theorems that need
it. Proof irrelevance, propositional and function extensionality, and
quotients are not used.

\begin{definition}[\lean{PSet}]
$\V:\Type_1$ is the inductive type with the constructor
$\mathsf{sup}:\Pi A{:}\Type_0.\ (A\to\V)\to\V$. For $x=\mathsf{sup}\,A\,f$ we write
$\bar x:=A$ for the \emph{index type} and $x_i:=f\,i$. Bisimulation
$x\approx y$ is defined by recursion:
$\forall i\,\exists j.\ x_i\approx y_j$ and $\forall j\,\exists i.\ x_i\approx y_j$.
Membership is $x\in y:=\exists j.\ x\approx y_j$.
\end{definition}

Bisimulation is an equivalence relation, membership respects it on both
sides, $x\approx y$ iff $x$ and $y$ have the same elements, and
$\in$-induction holds for every predicate on $\V$, whether or not it
respects $\approx$. We never quotient: a \emph{set} is an element of $\V$,
and the statements below are about $\approx$.

\paragraph{Operations.}
All of the following are terms, with the expected membership
conditions. The empty set; $\{A_i\mid i:\iota\}$ for a family indexed
by a small type; $\bigcup x$, and $\bigcup_{i:\iota}A_i$; pairs
$\{x,y\}$, the successor $x\cup\{x\}$, Kuratowski pairs $\langle x,y\rangle$, which
are injective up to $\approx$, and tuples; Separation
$\{z\in x\mid P\,z\}$ for an arbitrary $P:\V\to\Prop$, with index type the
subtype $\{i:\bar x\mid P\,x_i\}$; the power set, with index type
$\bar x\to\Prop$, which is small because $\Prop$ is impredicative;
$\omega$; and, for a proposition $P$ and $t:P\to\V$, the \emph{guard}
\[\guard tP:=\textstyle\bigcup_{h:P}t\,h,\qquad
z\in\guard tP\iff\exists h{:}P.\ z\in t\,h.\]
So $\V$ satisfies Zermelo set theory with Foundation, with Separation
for all propositional predicates. With excluded middle, ordinals
(transitive sets of transitive sets) are linearly ordered by $\in$. The
rank $\rk x:=\bigcup_i(\rk x_i\cup\{\rk x_i\})$ is a term, it is an
ordinal, and $\rk\eta\approx\eta$ for an ordinal $\eta$. The levels are given by
the term $V_x:=\bigcup_i\mathcal P(V_{x_i})$, defined for every set
$x$, with $y\in V_x$ iff $\rk y\in\rk x$. Note that $\rk$ and $V$ are
functions because they are defined by recursion on \emph{data}. What
$\V$ lacks is any way to obtain a set from a description of it.

\section{The materializing recursion}\label{sec:mat}

Fix a term $D:\V\to\V$ and a set $U:\V$. In the application $D(G)$
is $V_{\rk G+\omega}$ and $U$ is the domain of a functional relation.

\begin{definition}[\lean{Label}, \lean{Path}, \lean{step}, \lean{F}]\label{def:mat}
A \emph{label} is $\la$, or $\lb\,w$, or $\lc\,x$, for sets $w,x$; two
labels are equivalent if they have the same constructor and
bisimilar arguments. A \emph{path} is a list of labels, deepest first;
$[\,]$ is the root and $l::p$ is the child of $p$ with label $l$. Labels
and paths are types in $\Type_1$. For a relation $R$ on paths, a path $p$
and $\mathrm{rec}:\Pi c.\ R\,c\,p\to\V$ let
\begin{align*}
\mathrm{call}(l)&:=\guard{\mathrm{succ}\,(\mathrm{rec}\,(l::p)\,h)}{h:R\,(l::p)\,p},\\
G&:=\guard{\mathrm{rec}\,(\la::p)\,h}{h:R\,(\la::p)\,p},\\
\mathrm{step}_R\,p\,\mathrm{rec}&:=\mathrm{call}(\la)\ \cup
  \textstyle\bigcup_{i:\overline{D(G)}}\mathrm{call}(\lb\,D(G)_i)\ \cup
  \bigcup_{j:\bar U}\mathrm{call}(\lc\,U_j),
\end{align*}
and let $F_R:\Pi p.\ \Acc_R\,p\to\V$ be defined from $\mathrm{step}_R$ by the
eliminator of $\Acc$, so that
\[F_R\,p\,a=\mathrm{step}_R\,p\,(\lambda c\,h.\ F_R\,c\,a_{c,h}),\]
where $a_{c,h}:\Acc_R\,c$ is obtained from $a$ and $h:R\,c\,p$.
\end{definition}

The step calls the recursion at the child $\la$; then at the children
$\lb\,w$ for the elements $w$ of $D(G)$, where $G$ is the value of the
first call; and at the children $\lc\,x$ for $x\in U$. Every call is
guarded by the proposition that the child is below $p$, which the step
cannot prove and does not need to decide. The value is the union of the
successors of the values of the calls whose guard holds.

\begin{definition}[\lean{Coherent}]\label{def:coherent}
A \emph{target assignment} is a relation $\tau$ between paths and sets.
We say that $p$ \emph{has the target} $t$ if $\tau\,p\,t$. Its relation is
$R_\tau\,c\,p$: $c=l::p$ for some $l$, and $c$ has a target. Say that $G$
\emph{is the first value at $p$} if the elements of $G$ are the elements
of the targets of $\la::p$ (so $G$ is empty if $\la::p$ has no target),
and that a label $l$ is \emph{available for $G$} if $l=\la$, or $l$ is
equivalent to $\lb\,w$ with $w\in D(G)$, or to $\lc\,x$ with $x\in U$.
The assignment is \emph{coherent} if
\begin{enumerate}
\item the targets of a path are closed under $\approx$, and any two are
  bisimilar;
\item equivalent labels give children with the same targets;
\item if $l::p$ has the target $t$ and $p$ has the target $t'$ then
  $t\in t'$; and
\item if $p$ has the target $t$ and $G$ is the first value at $p$, then
  $x\in t$ iff $x\in t'\cup\{t'\}$ for some target $t'$ of a child $l::p$
  with $l$ available for $G$.
\end{enumerate}
\end{definition}

A coherent assignment is a proposition about a relation. No part of it
is data, and it may be defined using any quantifiers. Condition (3) does
not require the parent of a path with a target to have one.

\begin{lemma}[materialization; \lean{materialize}]\label{lem:mat}
If $\tau$ is coherent and $p$ has the target $t$, then $p$ is
$R_\tau$-accessible, and $F_{R_\tau}\,p\,a\approx t$ for every $a:\Acc_{R_\tau}\,p$.
\end{lemma}
\begin{proof}
By $\in$-induction on $t$, for all $p$. If $R_\tau\,c\,p$ then $c=l::p$ has a
target $t_c$, and $t_c\in t$ by (3); so the induction hypothesis applies
to every such $c$: it is accessible, and the recursion returns $t_c$ at
it. Hence $p$ is accessible, by the constructor. For the value, unfold
$F\,p\,a$ once. A guard at a child holds iff the child has a target, and
then the recursive call returns that target, up to $\approx$. So the set
$G$ of the step is the first value at $p$; the labels at which the step
makes calls are, up to equivalence, the labels available for $G$, and by
(2) equivalence does not matter; and the value of the step has as
elements the $x\in t'\cup\{t'\}$ for the targets $t'$ of the children with
available labels. By (4) this is $t$.
\end{proof}

The proof uses no excluded middle. The height of the recursion is not
bounded by anything computed from $p$: if the child $\la$ has value $G$,
the children $\lb\,w$ range over $D(G)$, which may be much larger than
anything seen so far, and so on below them.

\section{Replacement from uniform assignments}\label{sec:repl}

\begin{theorem}[\lean{replacement}]\label{thm:repl}
Let $C$ be a class of sets and $T_\eta$ a target assignment for each set
$\eta$, such that $T_\eta$ only depends on $\eta$ up to $\approx$ and, for $\eta$ in $C$,
$T_\eta$ is coherent for the parameter $U=s$ and gives the root the
target $\eta$. Let $\varphi$ be a relation that respects $\approx$ and is functional
on $s$, with all values in $C$. Then there is a set whose elements are
the $y$ with $\varphi(x,y)$ for some $x\in s$.
\end{theorem}
\begin{proof}
Let $\tau$ be the assignment that gives the path $r\mathbin{+\!\!+}[\lc\,x]$ the
target $t$ if $x\in s$, $\varphi(x,\eta)$, and $T_\eta\,r\,t$. It is \emph{one}
proposition, defined from $s$, $\varphi$ and $T$; no assignment is chosen for
each $x$. It is coherent, each condition reducing to the same condition
for some $T_\eta$, since the paths below $[\lc\,x]$ are in bijection with all
paths. The root has no target under $\tau$. Its target would be the union
of the $\eta\cup\{\eta\}$, and the existence of that set is the point at
issue. But every child of the root that has a target is accessible by
Lemma~\ref{lem:mat}, so the root is accessible by the constructor, and
$\theta:=F_{R_\tau}\,[\,]\,a$ is a term of type $\V$, because $s$ is a
variable of type $\V$. Unfolding the recursion once, the first value at
the root is empty, the children with a target are the $[\lc\,x]$ with
$x\in s$, and by the lemma the calls there return the values of $\varphi$. So
$y\in\theta$ iff $y\in\eta\cup\{\eta\}$ for some value $\eta$ of $\varphi$ on $s$, and the image is
$\{y\in\theta\mid\exists x\in s.\ \varphi(x,y)\}$.
\end{proof}

The assignments will come from a rule that acts at one node at a time.

\begin{lemma}[\lean{Rule.coherent}]\label{lem:rule}
Let $r(\eta,l,\xi)$ be a relation, the \emph{rule}, which determines $\xi$ up
to $\approx$ from $\eta$ and $l$ and respects $\approx$ and equivalence of labels. Let
$C$ be a class closed under $\approx$ such that for $\eta$ in $C$: if $r(\eta,l,\xi)$
then $\xi\in\eta$ and $\xi$ is in $C$; and, for $G$ with elements the
elements of the $\zeta$ with $r(\eta,\la,\zeta)$, we have $x\in\eta$ iff
$x\in\xi\cup\{\xi\}$ for some $l$ available for $G$ and $\xi$ with $r(\eta,l,\xi)$.
Define $T_\eta$ by recursion on the path: the root has the targets
bisimilar to $\eta$, and $l::p$ has the target $\xi$ if $p$ has a target $t$
with $r(t,l,\xi)$. Then $T_\eta$ is coherent for every $\eta$ in $C$.
\end{lemma}
\begin{proof}
Every target is in $C$, by induction on the path, and the conditions of
Definition~\ref{def:coherent} at $p$ are the hypotheses on the rule at the
target of $p$.
\end{proof}

\section{The definability rule}\label{sec:rule}

We need a rule whose class $C$ contains every ordinal. At a successor
$\zeta+1$ one child with target $\zeta$ suffices, and at $\omega$ the children can
be indexed by the natural numbers. At any other limit $\eta$ the children
must have targets cofinal in $\eta$, their labels must come from $D(G)$,
and $G$, being the target of a child, is below $\eta$. So the rule must
send sets of rank below $\eta$ cofinally into $\eta$, definably in $\eta$. It
cannot select one cofinal map, since $\ZF$ has no definable choice of a
cofinal map into a singular ordinal. The solution is not to select.

\begin{definition}[\lean{Unb}, \lean{Reach}, \lean{G$\nu$}]
Fix a relation $I(\eta,q,x,y)$ respecting $\approx$, read as ``$y$ is the value
at $x$ of the function defined over $V_\eta$ by the code $q$''. A pair
$(q,s)$ is \emph{unbounded at $\eta$} if $I(\eta,q,\cdot,\cdot)$ is functional on
$s$, every value at an $x\in s$ has rank in $\eta$, and for every $\zeta\in\eta$
some value at an $x\in s$ has rank at least $\zeta$. The ordinal $\eta$ is
\emph{reached from $\nu$} if some pair $(q,s)$ with $\rk q,\rk s\in\nu$ is
unbounded at $\eta$, and \emph{reachable} if it is reached from some
$\nu\in\eta$. For reachable $\eta$ let
$\nu(\eta):=\{\nu\in\eta\mid\eta\text{ is not reached from }\nu\}$, which is the
least ordinal from which $\eta$ is reached.
\end{definition}

The least $\nu$ is obtained by Separation and not by choosing: being
reached from $\nu$ is upwards closed in $\nu$, so the set of the $\nu$ that
fail is an ordinal, it is an element of $\eta$, and with excluded middle
$\eta$ is reached from it.

\begin{definition}[\lean{rule}]\label{def:rule}
With $D(G):=V_{\rk G+\omega}$, the rule $r(\eta,l,\xi)$ holds in the following
cases and no others.
\begin{enumerate}
\item $\eta\approx\zeta\cup\{\zeta\}$, $l=\la$, and $\xi\approx\zeta$.
\item $\eta\approx\omega$, and either $l=\la$ and $\xi\approx\emptyset$, or $l=\lb\,w$ with
  $\rk w\in\omega$ and $\xi\approx\rk w$.
\item $\eta$ is not a successor, not $\omega$, and reachable, and either
  $l=\la$ and $\xi\approx\nu(\eta)$, or $l=\lb\,w$ with $w\approx\langle q,s,x\rangle$, where
  $(q,s)$ is unbounded at $\eta$, $x\in s$, and $\xi$ is the rank of the value
  at $x$.
\end{enumerate}
\end{definition}

\begin{lemma}[\lean{worldly\_rule}]\label{lem:worldlyrule}
Assume excluded middle. Let $C$ be the class of ordinals $\eta$ such that
every ordinal $\mu\le\eta$ is $0$, a successor, $\omega$, or reachable. The rule
satisfies the hypotheses of Lemma~\ref{lem:rule} for $C$.
\end{lemma}
\begin{proof}
The cases are exclusive, and in each the target is determined by $\eta$
and the label: in (3) because a triple determines its components and
$(q,s)$ is functional on $s$. Every target is an element of $\eta$, and
$C$ is closed downwards. For the covering condition let $x\in\eta$. In (1),
$x\in\zeta\cup\{\zeta\}$ and $\la$ is available. In (2), $x$ is a natural number
$n$, the first value is $\emptyset$, and $n\in V_\omega=D(\emptyset)$. In (3),
$\eta$ is reached from $\nu(\eta)$, by a pair $(q,s)$; there is $x'\in s$ with a
value $y$ such that $x\in\rk y\cup\{\rk y\}$; the first value is $\nu(\eta)$; and
$\langle q,s,x'\rangle$ has rank below $\nu(\eta)+4$, since $q$, $s$ and $x'\in s$ have
rank below $\nu(\eta)$, so it is in $D(\nu(\eta))$. If $\eta$ is $0$ it has no
elements, and no child has a target.
\end{proof}

\begin{theorem}[\lean{dichotomy}]\label{thm:dich}
Assume excluded middle. Either every relation $\varphi$ that respects $\approx$ and
is functional on a set $s$ has an image in $\V$; or there is an ordinal
that is not $0$, not a successor, not $\omega$, and not reachable.
\end{theorem}
\begin{proof}
If every ordinal is $0$, a successor, $\omega$, or reachable, then $C$ is
the class of all ordinals. Given $s$ and $\varphi$, apply
Theorem~\ref{thm:repl} to the relation ``$\eta$ is the rank of the value of
$\varphi$ at $x$'', whose values are ordinals, with the assignments of
Lemmas~\ref{lem:rule} and~\ref{lem:worldlyrule}. Its image is a set of
ordinals, which is included in an ordinal $\theta$ by the proof of the
theorem, and the image of $\varphi$ is separated from $V_\theta$.
\end{proof}

The theorem holds for any $I$. The first alternative says more the
fewer ordinals are reachable, the second says more the more are; we now
choose $I$ so that the second gives a model of $\ZF$.

\section{$\ZF$ and its two models}\label{sec:zf}

\paragraph{Syntax and proofs (\lean{Fml}, \lean{Prf}, \lean{Con}).}
Formulas are built from $v_i\in v_j$, $v_i=v_j$, $\bot$, $\to$ and $\forall$,
with de Bruijn variables; $\neg$, $\wedge$, $\vee$, $\leftrightarrow$ and $\exists$ are
defined. The proof system is a Hilbert system: the propositional axioms
$K$ and $S$ and double negation elimination, modus ponens,
generalization, instantiation of $\forall$ by a variable,
$\forall(\varphi{\uparrow}\to\psi)\to\varphi\to\forall\psi$ where $\varphi{\uparrow}$ is $\varphi$ with its
variables shifted, reflexivity of equality, and substitutivity of
equality in atomic formulas. Nonlogical axioms are open formulas, their
free variables being parameters. $\Con(T)$ says that $T$ does not prove
$\bot$. Satisfaction $M\models\varphi[e]$ is defined for a class $M:\V\to\Prop$
and an environment $e:\N\to\V$, by recursion on $\varphi$ into $\Prop$, with
$\in$ and $\approx$ for the atoms and quantifiers ranging over $M$. Soundness
(\lean{soundness}): if every axiom of $T$ holds in $M$ under every
environment in $M$, so does every theorem; so a theory with a nonempty
class model is consistent. Double negation elimination is where
soundness uses excluded middle.

\paragraph{The axioms (\lean{ZF}).}
Extensionality; Foundation, as the existence of an $\in$-minimal element
of a nonempty set; Pairing, Union and Power set, in the form that a
set including the required one exists; Infinity, as a set containing an
empty set and closed under $y\mapsto y\cup\{y\}$; Separation for every formula
$\psi$, with arbitrary parameters,
$\exists y\,\forall z\,(z\in y\leftrightarrow z\in x\wedge\psi)$; and Replacement for every
formula $\psi$, with arbitrary parameters,
\[\forall x{\in}a\ \forall y\,y'\,(\psi(x,y)\to\psi(x,y')\to y=y')\ \to\ \exists b\ \forall y\,
  (\exists x{\in}a\ \psi(x,y)\to y\in b).\]
For each axiom a lemma states what its satisfaction in a class model
means in ordinary notation (\lean{sat\_ext}, \dots, \lean{sat\_repl}),
as a check that the de Bruijn indices say what is intended.

\paragraph{Models (\lean{ZFModel}).}
A class $M$ is a model of $\ZF$ if it is transitive, contains $\emptyset$ and
$\omega$, is closed under pairs, unions, power sets and subsets given by
arbitrary predicates, and for every formula $\psi$, parameters in $M$ and
$a$ in $M$ such that $\psi$ is functional on $a$ in the sense of $M$, some
$b$ in $M$ contains the values.

\begin{lemma}[\lean{V\_model}]
If every functional relation on a set has an image, the class of all
sets is a model of $\ZF$.
\end{lemma}

For the other alternative let $I(\eta,q,x,y)$ hold if $q\approx\langle\ulcorner\psi\urcorner,
\langle k,\langle e_0,\dots,e_{k-1}\rangle\rangle\rangle$ for a formula $\psi$ with free variables
below $k+2$ and a code $\ulcorner\psi\urcorner$ of it as a hereditarily finite set,
$x,y\in V_\eta$, and $V_\eta\models\psi[x,y,e_0,\dots,e_{k-1}]$ (\lean{ISat}). Codes and
tuples are injective up to $\approx$, so the code $q$ determines the
relation.

\begin{lemma}[\lean{Vl\_model}]\label{lem:vrho}
Assume excluded middle. If $\rho$ is an ordinal that is not $0$, not a
successor, not $\omega$, and not reachable for this $I$, then $V_\rho$ is a model
of $\ZF$.
\end{lemma}
\begin{proof}
By trichotomy $\omega\in\rho$ and $\rho$ is closed under successor. So finitely
many elements of $V_\rho$ have ranks below a common element of $\rho$, and
$V_\rho$ is closed under the operations, which raise ranks by a finite
amount. For Replacement let $\psi$ have its free variables below $k+2$, let
$e_0,\dots,e_{k-1}$ and $a$ be in $V_\rho$, with $\psi$ functional on $a$ in
$V_\rho$, and let $q$ be the code. The ranks of $q$ and $a$ are below some
$\nu\in\rho$: the code of $\psi$ has finite rank, and a tuple has rank a finite
amount above its components. The pair $(q,a)$ is functional on $a$, with
values in $V_\rho$. As $\rho$ is not reached from $\nu$, the pair is not
unbounded, so by excluded middle there is $\zeta\in\rho$ such that no value has
rank at least $\zeta$, and $b:=V_\zeta$ is as required.
\end{proof}

\begin{proof}[Proof of Theorem~\ref{thm:main}]
By Theorem~\ref{thm:dich} for this $I$, the two lemmas, and soundness.
\end{proof}

\paragraph{One model (\lean{HG}, \lean{hg\_model}).}
The two alternatives can be merged into a single model whose definition
involves neither a case distinction nor the ordinal $\rho$. Call an
ordinal \emph{good} if it is $0$, a successor, $\omega$, or reachable, and
\emph{hereditarily good} if every ordinal up to it is good; this is the
class $C$ of Lemma~\ref{lem:worldlyrule}. Let $H$ be the class of the
sets whose rank is hereditarily good. If every ordinal is good then $H$
is the class of all sets, and otherwise it is $V_\rho$ for the least
ordinal $\rho$ that is not good, but one does not need to know which.

\begin{theorem}[\lean{hg\_model}]
Assume excluded middle. Then $H$ is a model of $\ZF$.
\end{theorem}
\begin{proof}
$H$ is transitive, and it is closed under the set operations because
they raise ranks by at most one and the successor of a hereditarily good
ordinal is hereditarily good. For Replacement let $\psi$, with parameters
in $H$, be functional in $H$ on $a$ in $H$. The ranks of the values are
hereditarily good, so by Lemma~\ref{lem:worldlyrule} and
Theorem~\ref{thm:repl} they form a set, and their strict supremum $\theta$
exists; every ordinal below $\theta$ is hereditarily good. If $\theta$ is good,
then $V_\theta$ is in $H$ and contains the values. If $\theta$ were not good,
then, the hereditarily good ordinals being closed downwards, they would
be exactly the elements of $\theta$; so $H$ would be $V_\theta$, satisfaction in
$H$ would be satisfaction in $V_\theta$, and Lemma~\ref{lem:vrho} would
bound the values inside $V_\theta$. (In fact this case is contradictory: the
values would have ranks cofinal in $\theta$ and also bounded below it.)
\end{proof}
\section{The formalization}\label{sec:lean}

The development is a Lean~4 package with no dependencies. We did not
use Mathlib \cite{mathlib}, which has the sets-as-trees construction,
because its proofs use the axiom of choice freely and Lean's report of
the axioms used by a theorem would then say nothing. The files, in order
of dependency:
\begin{center}
\begin{tabular}{lrl}
\lean{PSet} & 146 & $\V$, $\approx$, $\in$, $\in$-induction, the basic operations, guards\\
\lean{Mat} & 149 & labels, paths, the recursion $F$, coherence, Lemma~\ref{lem:mat}\\
\lean{Repl} & 113 & the glued assignment, Theorem~\ref{thm:repl}\\
\lean{Rule} & 100 & Lemma~\ref{lem:rule}\\
\lean{Pair}, \lean{Ord}, \lean{VLevel} & 316 & pairs, ordinals, rank, power set, $V_x$, $\omega$, $D$\\
\lean{Worldly} & 291 & the rule, Lemma~\ref{lem:worldlyrule}, Theorem~\ref{thm:dich}\\
\lean{Fml}, \lean{Proof} & 283 & formulas, satisfaction, codes, the proof system, soundness\\
\lean{ZF}, \lean{ZFRead} & 275 & the axioms, models, readback of the axioms\\
\lean{FirstOrder} & 263 & $I$, Lemma~\ref{lem:vrho}, Theorem~\ref{thm:main}\\
\lean{Uniform} & 131 & the single model $H$\\
\lean{Stable}, \lean{Negative} & 123 & see the paragraph on excluded middle in Section~\ref{sec:discussion}
\end{tabular}
\end{center}
Excluded middle is never an axiom. It is a hypothesis
\lean{(em :\ $\forall$ p :\ Prop, p $\vee$ $\neg$p)} of the theorems that use it, which are:
trichotomy of ordinals and what depends on it (the characterization of
$V_x$ by rank, the least $\nu$), the case distinctions in
Lemma~\ref{lem:worldlyrule} and Theorem~\ref{thm:dich}, ``not unbounded
implies bounded'' in Lemma~\ref{lem:vrho}, and the classical connectives
in the semantics of formulas. Lemma~\ref{lem:mat},
Theorem~\ref{thm:repl} and Lemma~\ref{lem:rule} do not use it.

Lean reports that \lean{con\_ZF} depends on no axioms, and the source
asserts this with \lean{\#guard\_msgs}. Propositional extensionality is easy
to use by accident in Lean, and we record the three ways it entered: rewriting
with an equivalence (\lean{rw} with an \lean{Iff}); case analysis on an
equation between two distinct numerals other than $0$; and the library
lemmas about \lean{max} on natural numbers. We avoided them by
composing equivalences explicitly, by \lean{decide}, and by a
three-line definition of the maximum.

The universe level of $\V$ is a parameter throughout, and the main
theorem is proved at level $0$, where $\V:\Type_1$. The recursion $F$ is
an application of \lean{Acc.rec} with carrier the type of paths and
motive the constant $\V$, both in $\Type_1$.

\section{Discussion}\label{sec:discussion}

\paragraph{What the first alternative says.}
In the standard model of the type theory, where $\Type_0$ is
$V_\kappa$ for an inaccessible $\kappa$ and propositions are truth values,
$\V$ is $V_\kappa$ and the first alternative of
Theorem~\ref{thm:main} holds, as does the second. The first alternative
is Replacement for arbitrary propositional relations, the second-order
form, so where it holds the height of $\V$ is regular with respect to
everything the type theory can express. We do not know whether the type
theory proves it outright. The proof gives it from a weaker
hypothesis than the one we used: it suffices that for one relation $I$,
which may be any proposition and need not be first-order satisfaction,
every limit ordinal above $\omega$ is reachable.

\paragraph{More universes.}
With a third universe there is a type $\V_1$ of trees with index types
in $\Type_1$, which contains a set $\kappa_0$ of the ranks of the elements
of $\V$. We expect that the recursion of Section~\ref{sec:mat}, run on
paths of $\V$ with a relation that mentions a parameter from $\V_1$,
shows that $\kappa_0$ is a regular cardinal of $\V_1$, hence inaccessible
there, and so on up the hierarchy, which would give $\ZF$ with $n$
inaccessible cardinals from $n+2$ universes, matching the strength known
with choice \cite{werner97,carneiro19}. This is not checked.

\begin{question}
Is the type theory with $\omega$ universes and excluded middle, without
choice, equiconsistent with the same theory with choice, that is, with
$\ZFC$ with $n$ inaccessible cardinals for every $n$?
\end{question}

\paragraph{Small carriers.}
If $X$ is a type in $\Type_0$ and $R$ a relation on it, the accessible
part of $R$ is a well-founded relation on a small type. Its height $\alpha$
injects into the power set of $X$, by sending $\xi<\alpha$ to the set of
elements of height $\xi$, so it is below the Hartogs number of
$X\to\Prop$; and the Hartogs number of a small type $Y$ is an ordinal of
$\V$ given by a term (it is the set of the order types of the
well-orderings of subsets of $Y$, each of which is computed by recursion
on its own accessibility, on a small carrier). So a recursion
on $\Acc_R$ explores a tree that is bounded in advance, and the argument
of this paper has no analogue. The restriction is not artificial: the
well-founded recursion used in ordinary mathematics and programming is
mostly on small carriers.

\begin{question}\label{q:small}
Does $\ZF$ prove the consistency of the type theory with excluded
middle, without choice, in which $\Acc$ is restricted to carriers in
$\Type_0$?
\end{question}

\paragraph{Excluded middle.}
Since $\Con(\ZF)$ is a negation, it follows from the double negation of
excluded middle, which is equivalent to the double negation shift
$(\forall x.\,\neg\neg A\,x)\to\neg\neg\forall x.\,A\,x$ for predicates on arbitrary types
(\lean{con\_ZF\_of\_dns}). Whether the type theory proves $\Con(\ZF)$ with
no hypothesis is open, and we can say where the difficulty lies. Most
uses of excluded middle can be removed by a negative translation: read
bisimulation and membership with $\neg\neg\exists$ in place of $\exists$ at every
level of the trees (\lean{NEquiv}, \lean{NMem}); then every statement
about sets is stable under double negation, so it may be proved
classically, the semantics of formulas built from $\bot$, $\to$ and $\forall$
validates double negation elimination, and the case distinction of
Theorem~\ref{thm:dich} is only needed doubly negated. This includes
the theory of ordinals. The materialization lemma and
Theorem~\ref{thm:repl} are constructive as they stand.

The one statement that is not stable is accessibility. It is needed as
such, or at least doubly negated at the root, because the set that
bounds the ranks is the value of the recursion at the root. In
Lemma~\ref{lem:mat} accessibility is proved by $\in$-induction on the
target, from condition (3) of Definition~\ref{def:coherent} read with an
actual index: the target of a child is bisimilar to \emph{this} element
of the target of its parent. The targets produced by the rule are
specified by propositions (a least ordinal, the rank of the value of a
definable function), and after the translation a child's target is only
known not to fail to be an element. Induction along the translated
membership holds for stable predicates (\lean{nmem\_induction}) and not
for accessibility: passing from ``every child is not not accessible'' to
``not not every child is accessible'' is an instance of the double
negation shift. This is not an artefact. For the search relation of a
decidable predicate $P$ on $\N$, the root is accessible iff $\exists k.\,P\,k$
(\lean{exists\_of\_acc\_search}), so the principle that a relation which
is well-founded for stable predicates is accessible contains Markov's
principle, and its doubly negated form contains the double negation
shift for $\forall n\,\exists k.\,P\,n\,k$. We do not know whether the particular
instance needed here is provable. A proof of $\Con(\ZF)$ without any
hypothesis would seem to need either a way to make the descent of
targets honest, which a specification by propositions resists, or a
different source of large elimination.

\paragraph{Other systems.}
The argument uses nothing specific to Lean: an impredicative $\Prop$,
inductive types, and large elimination for accessibility are also
present in Coq and its successor Rocq, where the sets-as-trees have been
studied extensively \cite{werner97,barras10}. We expect the development
to port directly, with excluded middle as the only assumption.

\bibliographystyle{plain}
\bibliography{references}

\end{document}
